\section{The positing morphism, mechanized}

Positing's action on a valuation has, until now, lived only in prose: one
posit at claim point \(\varphi\) with charge \(q\) moves a valuation \(v\) to
\(v \vee_k \delta_{\varphi,q}\), the point mass \(q\) at \(\varphi\), \(\bot_k\)
elsewhere. This section gives that morphism its first mechanization. Every
result below is stated \emph{pointwise} --- \(\forall \psi, \dots\,\psi =
\dots\,\psi\) --- never as \texttt{Valuation} function equality, so no
\texttt{funext} is used anywhere in the underlying proofs.

\begin{proposition}[Pointwise valuation order]\label{thm:valuation-lek}\lean{Obsmetrologia.Valuation.leK}\leanok\uses{thm:valuation, thm:kle}
The pointwise knowledge order on valuations: \(v \le_k w\) iff
\(v\,\varphi \le_k w\,\varphi\) for every claim point \(\varphi\).
\end{proposition}

\begin{definition}[The point-mass valuation]\label{thm:delta}\lean{Obsmetrologia.Valuation.delta}\leanok\uses{thm:claim-point, thm:charge, thm:kBot-least}
\(\delta_{\varphi,q}\) is the valuation carrying charge \(q\) at \(\varphi\)
and \(\bot_k\) (absurd, told none) everywhere else.
\end{definition}

\begin{definition}[Positing, mechanized]\label{thm:posit}\lean{Obsmetrologia.Valuation.posit}\leanok\uses{thm:delta, thm:four-bilattice}
\texttt{posit}~\(\varphi\)~\(q\)~\(v\) is one knowledge-join update of \(v\)
by the point mass \(\delta_{\varphi,q}\): \(\texttt{posit}\,\varphi\,q\,v =
v \vee_k \delta_{\varphi,q}\), pointwise.
\end{definition}

\begin{theorem}[Abstaining is the identity]\label{thm:posit-kbot}\lean{Obsmetrologia.Valuation.posit_kBot}\leanok\uses{thm:posit, thm:kBot-least}
Positing \(\bot_k\) adds no knowledge anywhere: \texttt{posit}~\(\varphi\)~\(\bot_k\)~\(v\)
agrees with \(v\) at every claim point.
\end{theorem}

\begin{theorem}[Positing is inflationary]\label{thm:posit-inflationary}\lean{Obsmetrologia.Valuation.posit_inflationary}\leanok\uses{thm:posit, thm:valuation-lek}
A posit never retracts knowledge already held: \(v\,\psi \le_k
\texttt{posit}\,\varphi\,q\,v\,\psi\) for every \(\psi\).
\end{theorem}

\begin{theorem}[Same-cell posits merge by joining charges]\label{thm:posit-posit-same}\lean{Obsmetrologia.Valuation.posit_posit_same}\leanok\uses{thm:posit}
Two posits at the \emph{same} cell fold together rather than overwrite:
\(\texttt{posit}\,\varphi\,q_2\,(\texttt{posit}\,\varphi\,q_1\,v) =
\texttt{posit}\,\varphi\,(q_1 \vee_k q_2)\,v\), pointwise.
\end{theorem}

\begin{theorem}[Posits commute, unconditionally]\label{thm:posit-comm}\lean{Obsmetrologia.Valuation.posit_comm}\leanok\uses{thm:posit}
Two posits commute regardless of whether their cells coincide or are
related in any way: the fact is pure associativity/commutativity of
\(\vee_k\) on \texttt{Charge}, so no hypothesis on \(\varphi_1, \varphi_2\)
is needed.
\end{theorem}

\begin{theorem}[Positing is monotone]\label{thm:posit-monotone}\lean{Obsmetrologia.Valuation.posit_monotone}\leanok\uses{thm:posit, thm:valuation-lek}
\texttt{posit}~\(\varphi\)~\(q\) preserves the pointwise knowledge order:
\(v \le_k w \to \texttt{posit}\,\varphi\,q\,v \le_k \texttt{posit}\,\varphi\,q\,w\).
\end{theorem}

\begin{theorem}[Positing is idempotent]\label{thm:posit-idem}\lean{Obsmetrologia.Valuation.posit_idem}\leanok\uses{thm:posit-posit-same, thm:joinK-lattice}
Positing the same charge at the same cell twice equals positing it once,
from \cref{thm:posit-posit-same} and idempotence of \(\vee_k\).
\end{theorem}

\subsection*{The four Belnap speech-act classes}

\begin{definition}[Abstaining, asserting, denying, contesting]\label{thm:belnap-classes}\lean{Obsmetrologia.Valuation.abstaining, Obsmetrologia.Valuation.asserting, Obsmetrologia.Valuation.denying, Obsmetrologia.Valuation.contesting}\leanok\uses{thm:posit, thm:four-charges}
The four speech-act classes, named by the charge they posit: abstaining
posits \(\bot_k\) (told none, a no-op by \cref{thm:posit-kbot}); asserting
posits \(t\) (told true only); denying posits \(f\) (told false only);
contesting posits \(\top_k\) (told both).
\end{definition}

\begin{theorem}[Contesting absorbs, at its own cell]\label{thm:contest-absorbs}\lean{Obsmetrologia.Valuation.contest_absorbs}\leanok\uses{thm:belnap-classes, thm:posit}
Once \(\varphi\) has been contested, any further posit \emph{at \(\varphi\)}
leaves it at \(\top_k\): \(\texttt{posit}\,\varphi\,q\,(\texttt{contesting}\,\varphi\,v)\,\varphi
= (\texttt{contesting}\,\varphi\,v)\,\varphi\). This is stated \textbf{only}
at the cell \(\varphi\) itself --- off that cell, contesting leaves \(v\)
untouched, so no cell-free absorption law holds and none is claimed.
\end{theorem}

\subsection*{Naturality under converse-reindexing}

\begin{proposition}[Converse and its involution, at claim points]\label{thm:claim-converse-eq-iff}\lean{Obsmetrologia.Claim.converse_eq_iff}\leanok\uses{thm:predication-converse, thm:converse-involution}
\(\mathrm{converse}\,\psi = \varphi \leftrightarrow \psi =
\mathrm{converse}\,\varphi\), immediate from converse being an involution.
The fact naturality of \texttt{delta} under converse-reindexing turns on.
\end{proposition}

\begin{theorem}[\texttt{delta} is natural under converse-reindexing]\label{thm:delta-converse}\lean{Obsmetrologia.Valuation.delta_converse}\leanok\uses{thm:delta, thm:claim-converse-eq-iff, thm:reindex}
\(\delta_{\varphi,q}(\mathrm{converse}\,\psi) = \delta_{\mathrm{converse}\,\varphi,q}(\psi)\):
hitting the point mass after converse-ing the query point equals hitting the
converse-shifted point mass directly.
\end{theorem}

\begin{theorem}[\texttt{posit} is natural under converse-reindexing]\label{thm:posit-converse-natural}\lean{Obsmetrologia.Valuation.posit_converse_natural}\leanok\uses{thm:delta-converse, thm:posit, thm:reindex}
Reindexing a posited valuation along converse is the same as positing at the
converse cell into the already-reindexed valuation:
\(r^{*}(\texttt{posit}\,\varphi\,q\,v) = \texttt{posit}\,(\mathrm{converse}\,\varphi)\,q\,(r^{*}v)\)
for \(r = \texttt{reindex}\,\mathrm{converse}\), pointwise. The converse
generator commutes with positing.
\end{theorem}

The closure companion of this evidence-injection step --- how a
\emph{sequence} of posits settles to a fixed point under the modal closure
--- is \texttt{Language.Settling} territory (\cref{thm:settling}) and is not
addressed here. What this section supplies toward that closure's merge:
\cref{thm:posit-posit-same,thm:posit-comm,thm:posit-idem} together say that
any \emph{finite batch} of posits collapses, order-independently, to the
knowledge-join of the individually posited charges at each cell --- proved
here for two posits at a time; the finite-batch induction itself is carried
out in the next section.
